\section{讨论与未来工作}

本文方法的定位是小论文级别的探索性升级方案，而不是完整数据流通平台。其价值在于沿着 BLoP 已有的可信节点治理和 PBFT 确认主线，进一步区分不同风险等级下的执行路径：全节点模式继续承担高安全基线，轻量委员会模式承担高频审计下的成本优化。这样可以把现有研究中相对分散的组件串成一个闭环：PDP/PoR 提供完整性证明思想，区块链提供不可篡改记录和共识确认，纠删码提供异常后的恢复能力，可信评分和委员会机制降低中低风险事件的确认成本。

当前方案仍有明显边界。第一，本文没有实现真实 TPM、TEE 或 SGX 远程证明，可信评分只是对 BLoP 可信度量思想的软件化抽象。第二，当前校验流程以哈希抽样为主，不能等同于完整 PDP/PoR 密码学协议。第三，链上确认采用 PBFT 风格模型，尚未在 FISCO BCOS、Fabric 或 Tendermint 等真实联盟链平台上完成部署。第四，本文关注数据块和证据层面的完整性，不处理征信特征、隐私计算查询和联邦学习模型更新中的业务语义正确性。

后续工作可以从三个方向推进。第一，补充真实实验部分，在统一环境中比较全可信节点确认、Merkle 批量存证、可信委员会和 TrustFlowAudit 全量策略的链上记录数、确认延迟、存储冗余和恢复成功率。第二，将哈希抽样模块替换为 PDP/PoR 或零知识证明模块，提高安全表达强度。第三，将方法接入真实联盟链和隐私计算框架，进一步评估链下计算结果、授权日志和数据流通任务摘要的可审计性。
